\section{Conclusion}
\label{sec:conclusion}

\shvar{} combines Gaussian residual prediction with one shared encoder and a three-term objective: latent prediction, KL divergence and SIGReg. On Moving-MNIST, adding the residual to the shared recipe lowers the physical Energy Score in all five replications, and much of this gap is already present when real future embeddings are decoded, while the difference from the \emavar{} remains uncertain. On MPI3D, the residual gives no clear improvement over \shdet{}, and none of the ten variational-model runs beats its own persistence forecast on the gap-normalized latent score; the deterministic models cannot by construction. \shvar{} stores fewer parameters than the \emavar{} on both tasks. Its training compute per update was higher on Moving-MNIST, where the posterior needs a separate observation, and lower on MPI3D, where it reuses the target encoding. Natural next experiments would test prediction-scale control on MPI3D, score MPI3D forecasts on training attribute combinations to separate forecasting from compositional generalization, measure how often weak Moving-MNIST runs occur, and test whether digit identity can be preserved alongside the forecasting gains.
